% ===================================================================
\section{Method: Generalized Pair Specification}
\label{sec:method}

The Steersman's contrastive loss, as defined in
Paper~3~\cite{bielec2026bridge}, encodes the rhombic dodecahedron's
face-pair geometry through a fixed partition of channel indices into
co-planar and cross-planar sets. We generalize this mechanism by
replacing the hard-coded RD partition with a configurable pair
specification interface. Five experimental topologies test the
boundaries of this generalization.

\subsection{The Pair Specification Interface}
\label{sec:pairs}

The pair specification is a function
$\texttt{pairs}(n) \to \{(i_1, j_1), \ldots, (i_k, j_k)\}$ that
returns a set of \emph{co-axial} channel index pairs for a given
channel count~$n$. All remaining pairs of distinct indices are
classified as \emph{cross-axial}. The contrastive loss operates on
this partition identically to Paper~3: it rewards co-axial coupling
and penalizes cross-axial coupling. The Steersman's three control
laws (connectivity trending, directionality stagnation, stability
monitoring) operate without modification.

For $n = 6$, the RD specification is:
\begin{equation}
  \texttt{pairs}_{\text{RD}}(6) = \{(0,5),\; (1,4),\; (2,3)\}
  \label{eq:rd_pairs}
\end{equation}
These are the three face-pair axes of the rhombic dodecahedron. Each
pair shares a vanishing coordinate in the face-normal vectors,
producing three independent $2 \times 2$ blocks in the bridge.

For $n = 8$, the tesseract specification is:
\begin{equation}
  \texttt{pairs}_{\text{tess}}(8) = \{(0,1),\; (2,3),\; (4,5),\;
    (6,7)\}
  \label{eq:tess_pairs}
\end{equation}
These are the four co-axial vertex pairs of the tesseract, aligned
with the four coordinate axes of $\mathbb{R}^4$.

For $n = 4$, the octahedral specification is:
\begin{equation}
  \texttt{pairs}_{\text{oct}}(4) = \{(0,1),\; (2,3)\}
  \label{eq:oct_pairs}
\end{equation}
These are two of the octahedron's three vertex-opposition axes,
selected to match the channel count. The octahedron is the simplest
regular polytope admitting a pair decomposition and provides the
minimum-channel-count test of the geometric coherence hypothesis.

The interface imposes no constraint on which indices are paired.
Any partition of $n$~channels into $\lfloor n/2 \rfloor$ pairs
constitutes a valid specification. This generality enables the
non-geometric controls described below.

\subsection{Tesseract Topology ($n = 8$)}
\label{sec:tesseract_method}

The tesseract (4-cube, $\gamma_4$) has 16~vertices decomposing
into 4~co-axial pairs along the four coordinate axes of
$\mathbb{R}^4$. The pair specification
$\texttt{pairs}_{\text{tess}}(8)$ assigns one channel pair per
tesseract axis. If the Steersman generalizes beyond 3D, this
specification should produce four independent $2 \times 2$
blocks---a $4{+}4$ eigenvalue split.

\subsection{Octahedral Topology ($n = 4$)}
\label{sec:octahedral_method}

The octahedron has 6~vertices forming 3~co-axial pairs. At $n = 4$
channels, two of these axes are specified:
$\texttt{pairs}_{\text{oct}}(4) = \{(0,3),\; (1,2)\}$.
This represents the minimum geometry: only 2~co-axial pairs and
4~cross-axial pairs. The cross-axial/co-axial pair count ratio
($4/2 = 2$) is smaller than for the tesseract ($24/4 = 6$) or the
RD ($12/3 = 4$), providing the strongest co-axial signal per
channel pair.

\subsection{Wrong-Labels Topology ($n = 6$)}
\label{sec:wrong_labels_method}

To test whether the Steersman programs whatever topology it is told
or whether geometrically correct pair specifications are required,
we define a \emph{wrong-labels} specification: $n = 6$ channels
partitioned into three pairs by random assignment rather than RD
face-pair geometry. The Steersman's contrastive loss treats these
random pairs as co-axial and all other pairs as cross-axial. If the
Steersman is a general topology programmer accepting arbitrary
partitions, it should produce block-diagonal structure aligned with
the random pairs. If geometric coherence is required, the random
partition should fail.

\subsection{Circular Resonance Topology ($n = 6$)}
\label{sec:resonance_method}

The resonance topology derives pair specifications from
number-theoretic relationships among the tracked prime threads of the
companion research programme~\cite{bielec2026shape}. For $n = 6$, the
pairs are assigned by algebraic relationships: Sophie Germain pairing
($2(11) + 1 = 23 \to$ channels $(0, 2)$), consecutive-prime pairing
($29, 31 \to$ channels $(3, 4)$), and residue pairing (same residue
$\bmod\,6 \to$ channels $(1, 5)$). This specification tests whether
algebraic structure---without spatial embedding---can serve as a
topology specification.

\subsection{Emanation Architecture ($n = 6$)}
\label{sec:emanation_method}

The emanation architecture introduces hierarchical bridge structure.
Rather than independent bridge matrices per adapter (as in all prior
experiments), the emanation design maintains a single \emph{master
bridge} shared across all layers, with per-layer \emph{offset
matrices} that capture local variation:
\begin{equation}
  \mathbf{M}^{(l)} = \mathbf{M}_{\text{master}} +
    \boldsymbol{\Delta}^{(l)}
  \label{eq:emanation}
\end{equation}
where $l$ indexes the transformer layer and offsets are scaled by
$1/\sqrt{l}$. The Steersman monitors a \emph{coherence} metric---the
mean Frobenius norm of the offsets relative to the master---and
applies regularization pressure when individual layers fragment
excessively. The same RD-derived co-axial pairs are specified, but
the contrastive loss operates through the master matrix rather than
per-layer bridges directly.

\subsection{Training Setup}

All experiments use TinyLlama-1.1B-Chat with the same hyperparameters
as Paper~3: learning rate $2 \times 10^{-4}$, batch size~2, gradient
accumulation~8 (effective batch size~16), LoRA rank~24, identity
initialization, random seed~42, maximum sequence length~512,
alpaca-cleaned dataset, 10{,}000 training steps. The channel size
$s = r/n$ varies with~$n$: $s = 8$ ($n = 3$), $s = 6$ ($n = 4$),
$s = 4$ ($n = 6$), $s = 3$ ($n = 8$).

Table~\ref{tab:topology_specs} summarizes all pair specifications.

\begin{table}[H]
\centering
\caption{Pair specifications for all experimental topologies. Co-axial
  pairs receive positive contrastive gradient; all remaining
  off-diagonal index pairs are cross-axial and receive negative
  gradient. $|\mathcal{C}|$ = co-axial pair count,
  $|\mathcal{X}|$ = cross-axial pair count.}
\label{tab:topology_specs}
\small
\begin{tabular}{llclcc}
\toprule
Topology & $n$ & Co-axial pairs & Source & $|\mathcal{C}|$ & $|\mathcal{X}|$ \\
\midrule
RD (3D) & 6 & $(0,5), (1,4), (2,3)$ & Face opposition & 3 & 12 \\
Tesseract (4D) & 8 & $(0,1), (2,3), (4,5), (6,7)$ & Vertex opposition & 4 & 24 \\
Octahedral (3D) & 4 & $(0,1), (2,3)$ & Vertex opposition & 2 & 4 \\
Wrong-labels & 6 & Random partition & None (control) & 3 & 12 \\
Resonance & 6 & $(0,2), (3,4), (1,5)$ & Number-theoretic & 3 & 12 \\
\bottomrule
\end{tabular}
\end{table}
